%! TeX program = lualatex
% - - - - - - - - - - - - - - - - - - - - - - - - - - - - - - - - - - - - - - - 
% Dokumentinstillinger og viktige pakker
% - - - - - - - - - - - - - - - - - - - - - - - - - - - - - - - - - - - - - - - 
\documentclass[10pt,a4paper,norsk]{article}
\usepackage[top=2.5cm,right=2.5cm,bottom=2.5cm,left=2.5cm,footnotesep=1cm]{geometry}
\usepackage{etoolbox}             % If statements and boolean values
\usepackage{iftex}                % Conditional statements based on current latex engine
\usepackage[norsk]{babel}         % Support for other languages than english
\usepackage[T1]{fontenc}          % Font encoding
\usepackage[babel]{microtype}     % Improve text appearance in numerous ways
\usepackage{textcomp}             % Add extra symbols
\usepackage{mathtools}            % Improvements for math
\usepackage{amssymb}              % Extended math symbols


% - - - - - - - - - - - - - - - - - - - - - - - - - - - - - - - - - - - - - - - 
% Initialisering av boolske verdier
% - - - - - - - - - - - - - - - - - - - - - - - - - - - - - - - - - - - - - - - 
\providetoggle{margin}
\providetoggle{dark}
\providetoggle{endnotes}
\providetoggle{colored}


% - - - - - - - - - - - - - - - - - - - - - - - - - - - - - - - - - - - - - - - 
% Font oppsett
% - - - - - - - - - - - - - - - - - - - - - - - - - - - - - - - - - - - - - - - 
\usepackage[nomath,oldstylenums,light]{kpfonts} % Main (serif/sans-serif) font
\usepackage{newpxmath}            % Math font
%\usepackage[scaled]{beramono}     % Monospace font
\usepackage{fontspec}             % Custom fonts
\setmonofont{PxPlus ToshibaSat 8x16}[AutoFakeBold,AutoFakeSlant] % Monospace font


% - - - - - - - - - - - - - - - - - - - - - - - - - - - - - - - - - - - - - - - 
% Innlasting av pakker
% - - - - - - - - - - - - - - - - - - - - - - - - - - - - - - - - - - - - - - - 
\usepackage[T1]{url}\urlstyle{sf} % Clickable URLs
\usepackage{fullwidth}            % Full width content
\usepackage{hyperref}             % Clickable references
\usepackage{csquotes}             % Support quote marks in various languages
\usepackage{caption}              % Caption customization
\usepackage{ulem}                 % Double underline
\usepackage[makeroom]{cancel}     % Crossing out text
\usepackage{nicefrac}             % Nice looking fractals
\usepackage[compact]{titlesec}    % Title style
\usepackage{setspace}             % Set paragraph spacing
\usepackage{xcolor}               % Colors
\usepackage{enumitem}             % Lists
\usepackage{listings}             % Insert code
\usepackage{tikz}                 % Drawing of graphics
\usepackage{graphicx}             % Insert images
\usepackage{booktabs}             % Nicer looking tables
\usepackage{fancyhdr}             % Header / footer
\usepackage{lastpage}             % Last page number
\usepackage{marginnote}           % Margin notes
\usepackage{enotez}               % Endnotes
%\usepackage{verbatim}             % Comments
%\usapackage{siunitx}              % Formatting of units
%\usepackage{biblatex}             % Bibliography
%\usepackage{pdfpages}             % Insert pdf pages
%\usepackage{gensymb}              % Degree symbol
%\usepackage{cleverref}            % Cross-referencing


% - - - - - - - - - - - - - - - - - - - - - - - - - - - - - - - - - - - - - - - 
% Innlasting av TikZ-bibliotek
% - - - - - - - - - - - - - - - - - - - - - - - - - - - - - - - - - - - - - - - 
%\usetikzlibrary{}


% - - - - - - - - - - - - - - - - - - - - - - - - - - - - - - - - - - - - - - - 
% Sett institutt, emne, tittel og forfatter(e)
% - - - - - - - - - - - - - - - - - - - - - - - - - - - - - - - - - - - - - - - 
\newcommand{\fakultet}{Fakultet for teknologi, miljø- og samfunnsvitskap}
\newcommand{\institutt}{Institutt for datateknologi, elektroteknologi og realfag}
\newcommand{\emne}{DAT102}
\newcommand{\tittel}{Obligatorisk innlevering 3}
\newcommand{\forfatter}{Stephen Neba Fuh, \ Tord Johan Melheim, \\ Ebubekir Siddik Yuksel, \ Casper Eide Özdemir-Børretzen}
\newcommand{\dato}{xx. mars 2026}


% - - - - - - - - - - - - - - - - - - - - - - - - - - - - - - - - - - - - - - - 
% Sett boolske verdier
% - - - - - - - - - - - - - - - - - - - - - - - - - - - - - - - - - - - - - - - 
\settoggle{margin}  {true}  % Stor høyremargin med plass til notater
\settoggle{endnotes}{true}  % Konverter fotnoter til endenoter
\settoggle{dark}    {false} % Dark mode
\settoggle{colored} {true}  % Bruk farget tekst (i kildekode/lenker/etc.)


% - - - - - - - - - - - - - - - - - - - - - - - - - - - - - - - - - - - - - - - 
% Oppsett av fargepalett
% - - - - - - - - - - - - - - - - - - - - - - - - - - - - - - - - - - - - - - - 
% Apprentice fargepalett (https://github.com/romainl/Apprentice)
\definecolor{color0}{HTML}{1C1C1C}
\definecolor{color1}{HTML}{AF5F5F}
\definecolor{color2}{HTML}{5F875F}
\definecolor{color3}{HTML}{87875F}
\definecolor{color4}{HTML}{5F87AF}
\definecolor{color5}{HTML}{5F5F87}
\definecolor{color6}{HTML}{5F8787}
\definecolor{color7}{HTML}{6C6C6C}
\definecolor{color8}{HTML}{444444}
\definecolor{color9}{HTML}{FF8700}
\definecolor{color10}{HTML}{87AF87}
\definecolor{color11}{HTML}{FFFFAF}
\definecolor{color12}{HTML}{87AFD7}
\definecolor{color13}{HTML}{8787AF}
\definecolor{color14}{HTML}{5FAFAF}
\definecolor{color15}{HTML}{FFFFFF}
\definecolor{colorfg}{HTML}{BCBCBC}
\definecolor{colorbg}{HTML}{262626}
% Habamax fargepalett (https://github.com/habamax/vim-habamax)
\definecolor{brightCyan}{HTML}{87AFAF}
\definecolor{cyan}{HTML}{5F8787}
\definecolor{darkCyan}{HTML}{1F3F5F}
\definecolor{pink}{HTML}{D75F87}
\definecolor{red}{HTML}{D75F5F}
\definecolor{darkRed}{HTML}{AF5F5F}
\definecolor{brightGreen}{HTML}{5FF75F}
\definecolor{green}{HTML}{87D787}
\definecolor{darkGreen}{HTML}{5FAF5F}
\definecolor{blue}{HTML}{5fafd7}
\definecolor{darkBlue}{HTML}{5F87AF}
\definecolor{brightYellow}{HTML}{ffaf5f}
\definecolor{yellow}{HTML}{d7af87}
\definecolor{darkYellow}{HTML}{af875f}
\definecolor{brightMagenta}{HTML}{ff00af}
\definecolor{magenta}{HTML}{d787d7}
\definecolor{darkMagenta}{HTML}{af87af}
% Fargeoppsett
\colorlet{fgColor}{color0}
\colorlet{bgColor}{color15}
\colorlet{urlColor}{pink}
\newcommand{\hvlPath}{hvl-black.png}
\iftoggle{dark}{
    \colorlet{fgColor}{colorfg}
    \colorlet{bgColor}{color0}
    \renewcommand{\hvlPath}{hvl-white.png}
}{}
\colorlet{linkColor}{fgColor}
\pagecolor{bgColor}
\color{fgColor}
\newcommand{\colorIgnore}[1]{}
\newcommand{\textcolorIgnore}[2]{#2}
\iftoggle{colored}{}{
    \let\textcolor=\textcolorIgnore
    \let\color=\colorIgnore
}


% - - - - - - - - - - - - - - - - - - - - - - - - - - - - - - - - - - - - - - - 
% Egendefinerte kommandoer
% - - - - - - - - - - - - - - - - - - - - - - - - - - - - - - - - - - - - - - - 
\newcommand{\acshape}{\lsstyle\scshape\MakeLowercase}
\newcommand{\textac}[1]{\textls{\textsc{\MakeLowercase{#1}}}}
\newcommand{\mn}[1]{\marginnote{\mnText{#1}}}
\newcommand{\mnFig}[2]{\marginnote{#1\vspace{0.1cm}\par\mnText{#2}}}
\newcommand{\mnText}[1]{\textac{\textbf{\textcolor{color1}{Note:}}} #1}
\newcommand{\threestars}{\begin{center}\par * * * \end{center}}
\newcommand\N{\ensuremath{\mathbb{N}}}
\newcommand\R{\ensuremath{\mathbb{R}}}
\newcommand\Z{\ensuremath{\mathbb{Z}}}
\newcommand\Q{\ensuremath{\mathbb{Q}}}
\newcommand\C{\ensuremath{\mathbb{C}}}
\newcommand{\oppgave}[1]{\subsection*{Oppgave #1}}
\newcommand{\oppgaveDelStart}{\begin{enumerate}[leftmargin=*,itemsep=0.8cm,labelsep=1.5em,label=\alph*)]}
\newcommand{\oppgaveDelSlutt}{\end{enumerate}}
\newcommand{\oppgaveDel}[1]{\item[\textbf{#1})]}


% - - - - - - - - - - - - - - - - - - - - - - - - - - - - - - - - - - - - - - - 
% Oppsett av pakker/instillinger
% - - - - - - - - - - - - - - - - - - - - - - - - - - - - - - - - - - - - - - - 
% Sluttnote-oppsett
\microtypesetup{nopatch=footnote,letterspace=50}
\setenotez{list-name=Kilder,backref=true,mark-format=\bfseries,list-heading={\section*{#1}}}
\iftoggle{endnotes}{ \let\footnote=\endnote }{}
% Marginoppsett
\iftoggle{margin}{
    \geometry{right=5cm, marginparsep=0.8cm, marginparwidth=3.4cm}
    \fullwidthsetup{width=\linewidth+2.5cm}
}{}
\renewcommand*{\marginfont}{\footnotesize}
% Oppsett av topptekst og bunntekst
\pagestyle{fancy}\fancyhf{}
\fancyhead[L]{\normalsize\textac{\emne \ \ \textendash \ \ \tittel}}
\fancyfoot[C]{\footnotesize\textbf{\thepage} \tiny\textbf{ / } \footnotesize\textbf{\pageref*{LastPage}}}
\fancypagestyle{first}{\fancyhead[L,C,R]{}\renewcommand{\headrulewidth}{0pt}}
% Titteloppsett
%\titlespacing*{\section}{0cm}{0.8ex}{1.2ex}
\titlespacing*{\section} {0pt}{4ex plus 1ex minus .2ex}{0.8ex}
\titlespacing*{\subsection} {0pt}{3.25ex plus 1ex minus .2ex}{1.5ex plus .2ex}
\titlespacing*{\subsubsection}{0pt}{3.25ex plus 1ex minus .2ex}{1.5ex plus .2ex}
\titlespacing*{\paragraph} {0pt}{3.25ex plus 1ex minus .2ex}{1em}
\titlespacing*{\subparagraph} {\parindent}{3.25ex plus 1ex minus .2ex}{1em}
\titleformat{\section}{\normalfont\large\bfseries\acshape}{\thesection \ \ \textendash \ \ }{0em}{}
\titleformat{\subsection}{\normalfont\normalsize\bfseries}{\thesubsection.}{0.75em}{}
\titleformat{\subsubsection}{\normalfont\normalsize\itshape}{\thesubsubsection.}{0.75em}{}
% Oppsett av listings (visning av kildekode)
\lstset{literate={æ}{{\ae}}1{Æ}{{\AE}}1{ø}{{\o}}1{Ø}{{\O}}1{å}{{\aa}}1{Å}{{\AA}}1,
keywordstyle={\color{color1}}, 
commentstyle={\color{cyan}},
basicstyle={\scriptsize\ttfamily}, 
frame=none,rulecolor=\color{color7},
numbers=none,numberstyle={\scriptsize\ttfamily\color{color8}},
aboveskip=0.1cm,belowskip=1cm,columns=fixed,
breaklines=true,breakatwhitespace=false,keepspaces=true,
showspaces=false,showstringspaces=false,tabsize=4,
captionpos=b,inputencoding=utf8,extendedchars=true}
% Annet oppsett
\linespread{1.05}
\setlength{\parindent}{0pt}
\hypersetup{colorlinks=true, urlcolor=urlColor, filecolor=urlColor, linkcolor=linkColor, citecolor=linkColor, anchorcolor=linkColor, pdfauthor={\forfatter}, pdftitle={\emne - \tittel}}
\captionsetup{format=hang, font=footnotesize, labelfont=bf, labelsep=endash}
%\pretolerance = 3000
%\tolerance = 6000   \hbadness = \tolerance
%\setlength{\parskip}{0cm plus 1mm}


% - - - - - - - - - - - - - - - - - - - - - - - - - - - - - - - - - - - - - - - 
% Dokumentstart
% - - - - - - - - - - - - - - - - - - - - - - - - - - - - - - - - - - - - - - - 
\begin{document}


% Tittelside
% - - - - - - - - - - - - - - - - - - - - - - - - - - - - - - - - - - - - - - - 
\thispagestyle{first}
%\vspace*{0.8cm}
\par\includegraphics[scale=0.035]{\hvlPath}
\vspace{0.15cm}
%\par\normalsize \fakultet
\par\normalsize \institutt
\vspace{0.30cm}
\hrule
\vspace{0.30cm}
\par\Large \textbf{\emne \ \ $\mid$ \ \ \tittel}
\vspace{0.25cm}
\hrule
\vspace{0.35cm}
%\par\footnotesize{Gruppemedlemmer:}
%\vspace{0.1cm}
\par\normalsize \textbf{\forfatter}
\par\footnotesize \dato
\vspace*{0.5cm}
\normalsize

%\lstinputlisting[language=Java]{../src/TabellMengde.java}

\newcommand{\code}[1]{\footnotesize\texttt{#1}\normalsize}
\newcommand{\codeHl}[1]{\code{\textcolor{red}{#1}}}

% Hovedinnhold
% - - - - - - - - - - - - - - - - - - - - - - - - - - - - - - - - - - - - - - - 
\section*{Uke 10 \textendash \ Oppgave 4}
\oppgaveDelStart

\oppgaveDel{a} Se vedlagte filer \code{./src/TabellMengde.java} og \code{./test/TabellMengdeTest.java}.
\oppgaveDel{b} Se vedlagte filer \code{./src/LenketMengde.java} og \code{./test/LenketMengdeTest.java}.
    \oppgaveDel{c} Se vedlagte filer \code{./src/JavaSetToMengde.java} og \code{./test/JavaSetToMengdeTest.java}.
\oppgaveDel{d}

\begin{enumerate}[leftmargin=*,itemsep=1cm,labelsep=1.5em,label=\roman*.]
    \item
        \codeHl{boolean inneholder(T element)}
        \par $x \in A$
        \par For både \code{TabellMengde} og \code{LenketMengde} vil metoden gå gjennom alle elementer frem til det gitte elementet er funnet, eller konkludere med at det ikke finnes i mengden hvis det ikke ble funnet etter gjennomgang av alle elementer. I beste tilfelle er mengden tom og løkken der søket gjennomføres hoppes over, dette gir $O(1)$ i O-notasjon. I verste tilfelle har vi en mengde som inneholder elementer og en det letes etter et element som ikke ekisterer i mengden. Da må alle elementer sjekkes og dette gir $O(n)$ i O-notasjon, der $n$ er antall elementer i mengden.
    \item
        \codeHl{boolean erDelmengdeAv(MengdeADT<T> mengde)}
        \par $A \subseteq B \rightarrow \forall x \in A, x \in B$
        \par Metoden fungerer veldig likt for begge klassene, og i begge tilfeler sammenlignes to mengder. La oss kalle de mengde $A$ og mengde $B$. I metoden undersøkes det om $A$ er delmengde av $B$. Først blir det utført en sjekk for å forsikre at antall elementer i $A$ er like eller mindre enn antall elementer i $B$. Videre itereres det gjennom alle elementene i $A$ og det blir utført kall på \code{B.inneholder(A.element)} for alle elementer i $A$. Videre har vi to mulige utfall av iterasjonen. 1) Hvis \code{inneholder} returner \code{false} for et av elementene kan iterasjonen avsluttes tidlig og det kan konkluderes at $A$ \underline{ikke} er delmengde av $B$, så metoden returnerer \code{false}. 2) Hvis alle elementer har blitt sjekket uten funn av element fra $A$ som ikke er med i $B$ kan det konkluderes at $A$ er delmengde av $B$ og metoden returnerer \code{true}.
        \par I beste tilfelle har $A$ flere elementer enn $B$ og metoden avsluttes tidlig. Dette gir $O(1)$ i O-notasjon. I verste tilfelle er $A$ delmengde av $B$ og for alle elementer i $A$ må det gjøres kall på \code{B.inneholder(A.element)}. Gitt to mengder med $n$ antall elementer får vi da $n \cdot n$ antall ekvivalens operasjoner i \code{inneholder}. Med andre ord en worst case kjøretid uttrykt som $O(n^2)$ i O-notasjon.
    \item
        \codeHl{boolean erLik(MengdeADT<T> mengde)}
        \par $A = B$
        \par Metoden fungerer veldig likt for begge klassene og for to gitte mengder $A$ og $B$ gjennomføres det en binær operasjon som returnerer \code{true} hvis $A = B$ og \code{false} hvis $A \neq B$. Metoden gjennomfører først sjekk for å forsikre at $A$ og $B$ har likt antall elementer. Videre brukes den tidligere definerte metoden \code{A.erDelmengdeAv(B)} for å sjekke om alle elementer i $A$ også finnes i $B$. Dette gir samme O-notasjon for \code{erLik(...)} som for \code{erDelmengdeAv(...)} med best case $O(1)$ i tilfeller der $A$ og $B$ har ulikt antall elementer. Og worst case $O(n^2)$ i tilfeller der $A = B$.
    \item
        \codeHl{MengdeADT<T> union(MengdeADT<T> mengde)}
        \par $A \cup B$
        \par Metoden fungerer veldig likt for begge klassene \code{TabellMengde} og \code{LenketMengde}. Det opprettes en ny mengde og alle elementer fra $A$ legges direkte til den nye mengden (uten å gjennomføre sjekk om elementer allerede eksisterer, siden det ikke behøves). Videre prøves det å legge til alle elementer fra $B$ med metoden \code{A.leggTil(B.element)}, med innebygd sjekk \code{!A.inneholder(B.element)} for å unngå duplikater når nye elementer legges til. Dette gir en kjøretid på $n + n \cdot n$ der antall elementer i $A$ er $n$ og antall elementer i $B$ også er $n$ for enkelthets skyld (som tidligere). I O-notasjon blir dette $O(n^2)$ for både worst og best case.
    \item
        \codeHl{T fjern(T element)}
        \par $A - \{x\}$
        \par Denne motoden fungerer også svært likt for begge klassene. Gitt en mengde $A$ og et element $x$ som skal fjernes fra mengden, gjennomføres det en iterasjon gjennom mengden for å finne $x$. Dette kan gi to ulike resultater. 1) Hvis $x$ blir funnet fjernes $x$ fra mengden og $x$ returneres. 2) Hvis alle elementer i $A$ blir sjekket uten at $x$ ble funnet returneres \code{null}. I worst case finnes ikke $x$ i $A$ og alle elementer må sjekkes. Dette gir $O(n)$ i O-notasjon. I best case er $A$ en tom mengde og \code{null} returneres umiddelbart, som gir $O(1)$ i O-notasjon.
\end{enumerate}

\oppgaveDel{e}
    \par HashSet bruker en HashMap (også kalt dictionary i noen språk) datastruktur for å lagre verdier. Gitt en mengde $A$ og en verdi $x$ som vi ønsker å lagre, vil det gjennomføres en hash-funksjon på $x$ som gir en heltallsverdi $x_h$. Gitt at kapasasiteten til $A$ er $n$ gjennomføres det videre en modulus operasjon på $x_h$ slik at vi får en indeksverdi $x_i = x_h \mod n$ som kan brukes for å lagre $x$ i den underliggende tabellen. I noen tilfeller vil vi få verdier med samme indeksverdi, og dette kalles en kollisjon. Kollisjoner må løses slik at verdiene får tildelt unike indeksverdier.
    \par Dette er en meget simplifisert forklaring på hvordan en HashMap datastruktur fungerer. Men utifra dette ser vi at det blir mye raskere å sjekke om et element eksisterer eller ikke i en gitt mengde, sammenlignet med vår \code{inneholder(...)} metode fra \code{TabellMengde} og \code{LenketMengde}. En \code{inneholder(T element)} funksjon på et sett som bruker en underliggende HashMap datastruktør vil ha en kjøretid på $O(1)$ uttrykt i O-notasjon ettersom vi gjennom hash'ing får direkte tilgang til indeksverdien for verdien (pluss at det må tas hensyn til kollisjoner).

\oppgaveDel{f} Se vedlagte filer \code{./src/Person.java} og \code{src/Program.java}.

\oppgaveDelSlutt

\section*{Uke 11 \textendash \ Oppgave 4}
%\oppgaveDelStart
%\oppgaveDel{a} ...
%\oppgaveDelSlutt
...

% Sluttnoter
% - - - - - - - - - - - - - - - - - - - - - - - - - - - - - - - - - - - - - - - 
\iftoggle{endnotes}{ 
    \newpage
    \printendnotes
}{}


% - - - - - - - - - - - - - - - - - - - - - - - - - - - - - - - - - - - - - - - 
% Dokumentslutt
% - - - - - - - - - - - - - - - - - - - - - - - - - - - - - - - - - - - - - - - 
\end{document}

